% !TeX spellcheck = hu_HU
\documentclass[a4paper,12pt]{article}

% Magyar nyelvi beállítások
\usepackage[utf8]{inputenc}
\usepackage[T1]{fontenc}
\usepackage[magyar]{babel}

% Linkek
\usepackage[colorlinks, allcolors=blue]{hyperref}

% Szebb betűtípus
\usepackage{lmodern}

% Matematika
\usepackage{amsmath, amssymb, amsthm, amsfonts}

% Oldalmargók
\usepackage[margin=2.5cm]{geometry}

% Színes dobozok definíciókhoz, tételekhez stb.
\usepackage{tcolorbox}
\tcbuselibrary{theorems}

%coordinate
\usepackage{tikz}

% Tételstílusok
%no counter / number within=section
\newtcbtheorem[no counter]{definition}{Def.}%
{colback=blue!5,colframe=blue!50!black,fonttitle=\bfseries}{def}

\newtcbtheorem[no counter]{theorem}{Tétel}%
{colback=green!5,colframe=green!50!black,fonttitle=\bfseries}{th}

\newtcbtheorem[no counter]{example}{Példa}%
{colback=yellow!10,colframe=yellow!50!black,fonttitle=\bfseries}{ex}

\newtcbtheorem[no counter]{remark}{Megjegyzés}%
{colback=gray!10,colframe=gray!50!black,fonttitle=\bfseries}{rem}

\newtcbtheorem[no counter]{observation}{Megfigyelés}%
{colback=orange!5, colframe=orange!80!black, fonttitle=\bfseries}{obs}

\newtcbtheorem[no counter]{statement}{Allítás}%
{colback=orange!5, colframe=red!80!black, fonttitle=\bfseries}{sta}

% Fejezetcímekhez kicsit szebb stílus
\usepackage{titlesec}
\titleformat{\section}{\Large\bfseries}{\thesection.}{0.5em}{}
\titleformat{\subsection}{\large\bfseries}{\thesubsection.}{0.5em}{}
